\clearpage
\begin{proofbox}
\textbf{方法一：等价平方，只需一步消元（优先）。}
以顶点为原点建系，设 $P(x,y)$，$F(p/2,0)$，$\ell:x=-p/2$。
两边为非负距离，平方可逆：
\[
\begin{aligned}
PF=PH&\iff \sqrt{(x-p/2)^2+y^2}=|x+p/2|\\
&\iff (x-p/2)^2+y^2=(x+p/2)^2\\
&\iff \boxed{y^2=2px}.
\end{aligned}
\]
由方程 $x\ge0$，所以 $PF=PH=x+p/2\ge p/2>0$。
这同时验证了反向关系、比值分母不为零及最短焦半径。

\textbf{方法二：圆的相切判定（几何角度）。}
以 $P$ 为圆心、$PF$ 为半径作圆，圆当然经过 $F$。
直线与圆相切的充要条件是圆心到直线的距离等于半径，
因此它与 $\ell$ 相切 $\iff d(P,\ell)=PF$。
反过来，任意过 $F$ 且切 $\ell$ 的圆，其圆心都满足焦准等距。
所以圆心轨迹与本条抛物线完全相同；无须使用微积分或切线方程。
\end{proofbox}
\clearpage
\begin{proofbox}
\textbf{顶点为什么必须先定位？}
顶点到焦点和准线等距，且在焦点与准线垂足连线上，故恰为该线段的中点。
\begin{center}\begin{tikzpicture}[font=\small,>=stealth]
\draw[black!30,rounded corners=2pt] (0,0) rectangle (14.4,8.2);
\node[anchor=west,font=\small\bfseries] at (.5,7.72) {同一焦点 $F(3,0)$：准线改变，顶点和开口都会改变};
\begin{scope}[shift={(3.2,4)},scale=.4]
\draw[->,black!40] (-4,0)--(7,0) node[right] {$x$};
\draw[blue,very thick,domain=-4:4,samples=121,variable=\t] plot ({\t*\t/12},\t);
\draw[violet,thick] (-3,-4)--(-3,4) node[above] {$\ell$};\node[fill=black,circle,inner sep=1.8pt] at (3,0) {};\node[below left=5pt] at (3,0) {$F$};
\node[fill=black,circle,inner sep=1.8pt] at (0,0) {};
\end{scope}
\node at (3.8000000000000003,1.3) {$V=(0,0),\ y^2=12x$};
\begin{scope}[shift={(9.6,4)},scale=.4]
\draw[->,black!40] (-4,0)--(7,0) node[right] {$x$};
\draw[blue,very thick,domain=-4:4,samples=121,variable=\t] plot ({4-\t*\t/4},\t);
\draw[violet,thick] (5,-4)--(5,4) node[above] {$\ell$};\node[fill=black,circle,inner sep=1.8pt] at (3,0) {};\node[below left=5pt] at (3,0) {$F$};
\node[fill=black,circle,inner sep=1.8pt] at (4,0) {};
\end{scope}
\node at (10.2,1.3) {$V=(4,0),\ y^2=-4(x-4)$};

\node[anchor=west] at (.7,.55) {同尺度；左图准线 $x=-3$，右图准线 $x=5$，不能只凭焦点正负写标准式};
\end{tikzpicture}\end{center}

右图焦点也为 $(3,0)$，但准线为 $x=5$，顶点为 $(4,0)$，开口向左。
直接列式 $(x-3)^2+y^2=(x-5)^2$，得 $y^2=-4(x-4)$。
因此“焦点横坐标为正”不能替代“顶点在原点”这个标准位置条件。
一般平移后，若 $F=(h+p/2,k)$、$\ell:x=h-p/2$，方程为 $(y-k)^2=2p(x-h)$；
其他开口方向由反射或交换坐标得到。
\end{proofbox}

% ===== 以下为 C041 分层证明；原先两法与 TikZ 代码一字不改 =====
\clearpage
\begin{proofbox}
\textbf{方法三：直角三角形 + 平方差（最低代数门槛）。}
\begin{center}\begin{tikzpicture}[font=\small,>=stealth]
\draw[black!30,rounded corners=2pt] (0,0) rectangle (14.4,8.2);
\node[anchor=west,font=\small\bfseries] at (.5,7.72) {同取 $p=4$：勾股中的 $PQ$ 必须是非负长度};

\node at (3.6,6.7) {$0<x<p/2$：取 $P=(1/2,2)$};
\node at (10.8,6.7) {$x>p/2$：取 $P=(4,4\sqrt2)$};
\foreach \pos/\px/\py in {3.6/.5/2,10.8/4/5.6568542495}{
\begin{scope}[shift={(\pos,2.55)},scale=.55]
\draw[violet,thick] (-2,-.5)--(-2,6.5);
\draw[black!50,dashed] (2,0)--(2,6.5);
\draw[black!60,thick] (\px,\py)--(2,\py)--(2,0);
\draw[orange!85!black,very thick] (\px,\py)--(2,0);
\draw[violet,very thick] (-2,\py)--(\px,\py);
\ifdim\px pt>2pt \def\dx{.25}\else\def\dx{-.25}\fi
\draw[black!50] ({2+\dx},\py)--({2+\dx},{\py-.25})--(2,{\py-.25});
\fill (-2,\py) circle (2pt) node[left=6pt] {$H$};
\fill (\px,\py) circle (2pt) node[above=7pt] {$P$};
\fill (2,\py) circle (2pt) node[below right=7pt] {$Q$};
\fill (2,0) circle (2pt) node[below=7pt] {$F$};
\end{scope}}
\node at (3.6,1.45) {$PH=x+2,\quad PQ=2-x$};
\node at (10.8,1.45) {$PH=x+2,\quad PQ=x-2$};

\node at (7.2,.55) {两图同尺度；均有 $PF^2=PQ^2+y^2$，故 $y^2=(x+2)^2-(x-2)^2=8x$};
\end{tikzpicture}\end{center}

先看图中 $P$ 在 $x=p/2$ 两侧的两种情况。统一写法是
\[
PQ=|x-p/2|,\quad QF=|y|,\quad PH=|x+p/2|.
\]
勾股给出 $PF^2=PQ^2+QF^2$；即使 $P,Q,F$ 有重合点，这个距离平方式仍成立。
\end{proofbox}
\clearpage
\begin{proofbox}
\textbf{方法三续：先写平方，最后再判断符号。}
设 $P=(x,y)$、$F=(p/2,0)$、$H=(-p/2,y)$、$Q=(p/2,y)$。
由勾股距离式和焦准等距，
\[
\begin{aligned}
PF^2&=(x-p/2)^2+y^2,\\
PH^2&=(x+p/2)^2,\\
PF=PH&\iff y^2=(x+p/2)^2-(x-p/2)^2\\
&=\bigl[(x+p/2)-(x-p/2)\bigr]
  \bigl[(x+p/2)+(x-p/2)\bigr]\\
&=\boxed{2px}.
\end{aligned}
\]
这里是在对\textbf{代数式}作平方差，不能一开始就把 $x-p/2$ 当作长度。
由方程推出 $x\ge0$ 后，才有 $PH=x+p/2$；
但 $PQ$ 仍须分为 $p/2-x$ 或 $x-p/2$。
例如左图 $PH=5/2$、$PQ=3/2$，两长度之差是 $1=2x$、之和是 $4=p$；
右图才是两长度之差为 $p$、之和为 $2x$。两种位置的平方差都等于 $2px$。

\medskip
\textbf{考场起手式：}“先勾股，写距离平方，再用平方差”。
两边非负，平方可逆，故同时证明整个等距轨迹与标准方程等价。
\end{proofbox}
\clearpage
\begin{proofbox}
\textbf{方法四：三角形不等式，不建坐标就找最短焦半径。}
\begin{center}\begin{tikzpicture}[font=\small,>=stealth]
\draw[black!30,rounded corners=2pt] (0,0) rectangle (14.4,8.2);
\node[anchor=west,font=\small\bfseries] at (.5,7.72) {最短焦半径：两处不等式同时取等号，点 $P$ 才能落在顶点};

\node at (3.6,6.7) {一般点：$P=(2,4)$};
\node at (10.8,6.7) {等号点：$P=V=(0,0)$};
\foreach \pos in {3.6,10.8}{
\begin{scope}[shift={(\pos,3)},scale=.65]
\draw[violet,thick] (-2,-.8)--(-2,4.8);
\draw[black!55,dashed] (-2,0)--(2,0);
\fill (2,0) circle (2pt) node[below right=7pt] {$F$};
\end{scope}}
\begin{scope}[shift={(3.6,3)},scale=.65]
\draw[black!60,thick] (-2,4)--(2,0);
\draw[orange!85!black,very thick] (2,0)--(2,4);
\draw[violet,very thick] (-2,4)--(2,4);
\draw[black!55] (-2,.3)--(-1.7,.3)--(-1.7,0);
\fill (-2,0) circle (2pt) node[below left=7pt] {$D$};
\fill (-2,4) circle (2pt) node[above left=7pt] {$H$};
\fill (2,4) circle (2pt) node[above right=7pt] {$P$};
\end{scope}
\begin{scope}[shift={(10.8,3)},scale=.65]
\draw[violet,very thick] (-2,0)--(0,0);
\draw[orange!85!black,very thick] (0,0)--(2,0);
\fill (-2,0) circle (2pt) node[below left=7pt] {$H=D$};
\fill (0,0) circle (2pt) node[above=7pt] {$P=V$};
\end{scope}
\node at (3.6,1.55) {$FD=4<FH=4\sqrt2<2PF=8$};
\node at (10.8,1.55) {$FD=FH=2VF=4$};

\node at (7.2,.55) {两图同尺度，$p=4$；$H=D$ 且 $P$ 是 $FH$ 中点时，$PF=p/2$};
\end{tikzpicture}\end{center}

令 $D,H$ 分别是 $F,P$ 到准线的垂足，$FD=p$。只用两条初中几何事实：
\[
p=FD\ \underset{\text{垂线段最短}}{\le}\ FH
\ \underset{\text{两边和不小于第三边}}{\le}\ FP+PH=2FP.
\]
所以 $\boxed{PF\ge p/2}$。图中左侧两处都严格小于，右侧才同时取等号。

\textbf{等号不能漏查：}若 $PF=p/2$，则 $FD=FH$ 迫使 $H=D$；
$FH=FP+PH$ 迫使 $P$ 在线段 $FH$ 上。再由 $FP=PH$，
$P$ 只能是 $FD$ 中点 $V$。反过来 $VF=VD=p/2$，故等号恰在顶点。

此法专门证明最短焦半径，无须先求抛物线方程。
\end{proofbox}
\clearpage
\begin{proofbox}
\textbf{方法五：垂直平分线构造（纯几何识别轨迹）。}
已知 $F\notin\ell$。准线上任选 $H$；过 $H$ 作 $m\perp\ell$，
再作线段 $FH$ 的垂直平分线 $n$，令 $P=m\cap n$。

\textbf{正向：} $P\in n$ 说明 $PF=PH$（垂直平分线性质）；
$P\in m$ 说明 $PH=d(P,\ell)$（垂足是 $H$）；因此 $P$ 在抛物线上。

\textbf{反向：}若 $P$ 原本满足 $PF=d(P,\ell)$，令 $H$ 为垂足，
则 $PF=PH$，所以 $P$ 必在 $FH$ 的垂直平分线上，并且必在过 $H$ 的准线垂线上。
因此上述作图\textbf{既没有漏点，也没有多点}。

\textbf{为什么交点始终存在而且唯一？}
$F\notin\ell$ 使得线段 $FH$ 在准线的法向上有不为零的分量 $p$。
其垂直平分线 $n$ 不会与 $m$ 平行，故二者恰有一个交点。

\textbf{补一行代数，和标准方程接上：}
令 $F=(p/2,0)$、$H=(-p/2,t)$、$P=(x,t)$。
因 $PF=PH$，有
\[
\left(x-\frac p2\right)^2+t^2
=\left(x+\frac p2\right)^2
\quad\Longrightarrow\quad
\boxed{x=\frac{t^2}{2p}}.
\]
所以这种折纸构造真的描出了 $y^2=2px$，不是只得到一条“差不多”的曲线。
\end{proofbox}
\clearpage
\begin{proofbox}
\textbf{方法六：任意焦点与竖直准线直接列式（考场最快解方程法）。}
给定 $F(u,v)$、$\ell:x=c$，只要求 $u\ne c$。
设 $P(x,y)$，\textbf{不必先平移原点}，直接写
\[
\begin{aligned}
PF^2&=d(P,\ell)^2\\
\Longleftrightarrow\quad (x-u)^2+(y-v)^2&=(x-c)^2\\
\Longleftrightarrow\quad
&\boxed{(y-v)^2=(u-c)(2x-u-c)}.
\end{aligned}
\]
设 $h=(u+c)/2$，就得到
\[
\boxed{(y-v)^2=2(u-c)(x-h)}.
\]
这里 $V=(h,v)$ 是焦点与准线垂足的中点。
当 $u>c$ 时向右开口，$p=u-c$；当 $u<c$ 时向左开口，$p=c-u$。
这两个分支\textbf{不要把 $p>0$ 误写成有符号的 $u-c$}。

\textbf{若准线是水平的 $\ell:y=c$：}
同理，对焦点 $F(u,v)$ 有
\[
\boxed{(x-u)^2=(v-c)(2y-v-c)}.
\]
\textbf{选法建议：}已知焦点和水平、竖直准线，且只求方程，
“直接平方，两个平方差相消”通常最省事；
要讲清图像、开口、顶点时，再先用中点定位。
\end{proofbox}
\clearpage
\begin{proofbox}
\textbf{方法七：准线是斜线时的通法（进阶，不要求死记）。}
设 $\ell:ax+by+c=0$，其中 $a^2+b^2>0$，焦点 $F=(u,v)$ 且 $au+bv+c\ne0$。
任取 $P=(x,y)$，按两种距离的定义，直接有
\[
\boxed{(x-u)^2+(y-v)^2
=\frac{(ax+by+c)^2}{a^2+b^2}}.
\tag{1}
\]
这是\textbf{不旋转坐标}也能写出的轨迹方程：把点到线公式代进焦准等距，再平方即可。
两边原来都是非负的距离，因此平方是等价变形，反向没有增根。

\medskip
若想看清它为什么仍是一条抛物线，可按以下三步旋转思路：
\begin{enumerate}
\item 计算焦点在准线上的垂足
\[
D=F-\frac{au+bv+c}{a^2+b^2}(a,b),\qquad
p=\frac{|au+bv+c|}{\sqrt{a^2+b^2}}>0.
\]
\item 取顶点 $V=(F+D)/2$；取从准线指向焦点的单位法向量
\[
\boldsymbol n=\frac{\operatorname{sgn}(au+bv+c)}{\sqrt{a^2+b^2}}(a,b),
\qquad \boldsymbol t=(-n_y,n_x).
\]
\item 令 $s=(P-V)\cdot\boldsymbol n$（沿开口方向），
$r=(P-V)\cdot\boldsymbol t$（垂直于轴）。
在新的直角坐标里，焦点为 $(p/2,0)$、准线为 $s=-p/2$，
于是 $\boxed{r^2=2ps}$。
\end{enumerate}
\textbf{结论：}一般位置的抛物线并非一定能\emph{仅靠平移}写成 $y^2=2px$；
若准线倾斜，通常还需要旋转坐标。式 (1) 是更通用的速写法。
\end{proofbox}
